%=====================================================================
\chapter{Recent Trends and Applications}\label{ch:trends}
%=====================================================================
\locator{6}{Part VI --- Learning, Language and Responsible AI}

\begin{outcomes}
By the end of this chapter you will be able to:
\begin{tight}
  \item \textbf{Describe} what a large language model is, in terms this course has
        already established.
  \item \textbf{Explain} the structure of a tool-using agent, and \textbf{relate}
        it to the agent architectures of \chref{agents}.
  \item \textbf{Identify} which parts of a modern AI system are symbolic and which
        are learned.
  \item \textbf{Assess} a claim about an AI system using \chref{introduction}'s
        checklist and \chref{classic}'s framework.
  \item \textbf{Say} what this course's material contributes to systems built
        today.
\end{tight}
\end{outcomes}

\begin{prereq}
\chref{agents} (agent architectures), \chref{expert} (rule engines, which reappear
inside agent loops), \chref{nlp} (n-grams and the Eliza effect) and
\chref{learning} (what learning is). This chapter is a signpost: it names what is
current and connects it to what you have built.
\end{prereq}

\begin{keyterms}
large language model $\bullet$ token $\bullet$ context window $\bullet$ prompt
$\bullet$ tool use $\bullet$ function calling $\bullet$ agent loop $\bullet$
retrieval-augmented generation $\bullet$ neuro-symbolic $\bullet$ hallucination
$\bullet$ grounding $\bullet$ benchmark
\end{keyterms}

\section{Where the Field Is Now}

This course has taught the symbolic half of artificial intelligence, and it has been
honest about the fact that the systems in the news are not built that way. This
chapter says what they are, in the vocabulary you now have, and what your material
contributes to them.

It is a signpost rather than a unit. A serious treatment of the current wave needs
the statistics of \emph{Machine Learning} and the architectures of \emph{Data
Science}; what can be done here is to place it.

\section{Large Language Models, Briefly}

A \term{large language model} predicts the next \term{token} --- roughly, the next
word-piece --- given the tokens before it. That is the n-gram model of
\chref{nlp}, with two differences that turn out to matter enormously.

\textbf{The window is much larger, and it is not fixed by counting.} An n-gram
model conditions on the previous two or three words because counting longer
sequences is hopeless: the data is too sparse. A language model conditions on
thousands of tokens, because it does not count --- it learns a function, and the
function generalises across contexts that never co-occurred in training. That
directly addresses the limitation \chref{nlp} said was unfixable by counting.

\textbf{The scale is different in kind.} Hundreds of billions of parameters,
trained on a substantial fraction of published text. The empirical finding of the
last decade is that capabilities appear as this scales which nobody designed and
few predicted.

\begin{alertbox}[What it still is, and what follows]
It is a next-token predictor. Everything else --- answering questions, writing code,
apparently reasoning --- is a consequence of that objective being pursued at scale
over text that contains people doing those things.

Three consequences follow directly, and each one is visible in this course's terms.

\emph{There is no representation to inspect.} Nothing corresponds to
\chref{kr}'s semantic network or \chref{bayesnets}' conditional probability table.
So \chref{introduction}'s third test scores zero, and there is no artefact for a
domain expert to correct.

\emph{There is no proof and no trace.} A language model can \emph{produce} text
resembling \chref{expert}'s explanation chain, and that text is generated by the
same process as the answer --- it is not a record of how the answer was reached. The
distinction is exactly \chref{propositional}'s $\models$ against $\vdash$: a
plausible-looking derivation is not a derivation.

\emph{Fluency is not evidence.} This is \chref{nlp}'s Eliza effect at scale, and
your defence against it is \chref{introduction}'s checklist rather than your
impression of the transcript.
\end{alertbox}

\term{Hallucination} --- fluent, confident, false output --- is not a defect to be
patched but the objective working as specified. The model was asked for likely text,
and false text is often extremely likely. In \chref{classic}'s framework, its failure
mode is Eliza's: \emph{none detectable by the system}, which is the worst entry in
that column.

\section{Tool-Using Agents}

The interesting engineering development is not the model but the loop around it.

\dsfig{%
\begin{tikzpicture}[font=\scriptsize,
  b/.style={rounded corners=3pt, draw=#1, line width=1.1pt, fill=#1!10,
            text=#1!55!black, font=\scriptsize\bfseries, align=center,
            minimum width=2.7cm, minimum height=0.85cm},
  t/.style={rounded corners=3pt, draw=goldhl, line width=1pt, fill=goldhl!10,
            text=goldhl!50!black, font=\tiny\bfseries, align=center,
            minimum width=2.5cm, minimum height=0.6cm},
  a/.style={-{Stealth[length=2.2mm]}, neutral!75, line width=1pt}]

\node[b=primary]   (goal) at (-4.6,1.6) {goal};
\node[b=secondary] (dec)  at (0,1.6)    {decide next step};
\node[b=greenhl]   (act)  at (0,-0.3)   {call a tool};
\node[b=accent]    (obs)  at (-4.6,-0.3){observe result};
\node[b=purplehl]  (done) at (4.6,1.6)  {answer};

\node[t] (t1) at (4.6,-0.1)  {estimate(task)};
\node[t] (t2) at (4.6,-0.85) {dependencies(task)};
\node[t] (t3) at (4.6,-1.6)  {capacity(sprint)};

\draw[a] (goal) -- (dec);
\draw[a] (dec) -- node[right, font=\tiny, text=neutral] {not enough yet} (act);
\draw[a] (act) -- (obs);
\draw[a] (obs) .. controls +(0,1.0) and +(-1.6,0) .. (dec);
\draw[a] (dec) -- node[above, font=\tiny, text=neutral] {enough} (done);
\draw[a, goldhl] (act) -- (t1.west);
\draw[a, goldhl] (act) -- (t2.west);
\draw[a, goldhl] (act) -- (t3.west);

\node[rounded corners=3pt, fill=lightbg, draw=primary!30, line width=0.9pt,
      text=primary!70!black, font=\scriptsize, text width=13.4cm,
      align=flush left, inner sep=6pt] at (0,-3.0)
  {\textbf{Compare \dsref{agent-kinds}.} This is the model-based, goal-based agent
   of \chref{agents}: it keeps state (what the tools have returned), it holds a
   goal, and it selects actions to reach it. The architecture is thirty years old
   and the vocabulary of Part~I describes it exactly.

   \smallskip
   What is new is \emph{what decides the next step}. The classical answer was
   search (\chref{informed}) or a rule base (\chref{expert}); the current answer is
   a language model. Everything else --- the loop, the step budget, the tools, the
   termination test --- is ordinary engineering, and most of what makes such a
   system work or fail lives there rather than in the model.};
\end{tikzpicture}}%
{A tool-using agent. The loop is \chref{agents}' goal-based agent; only the
decision procedure at its centre has changed.}{agent-tools}

The chapter's lab builds this loop with a rule engine at its centre instead of a
language model --- deterministic, offline, and fully traceable. That substitution is
the point of the exercise: it isolates what the loop contributes from what the model
contributes, and the loop turns out to contribute most of the structure.

\term{Retrieval-augmented generation} is the same idea with one tool: search a
document collection, put what you find into the prompt, and generate from that. It
is the standard response to hallucination and it is a partial one --- the retrieved
text \term{grounds} the answer, and nothing forces the answer to respect it.

\section{Neuro-Symbolic Systems}

The current wave is strong exactly where this course's methods are weak, and weak
exactly where they are strong.

\smallskip
\noindent\begin{longtable}{@{}L{3.9cm}L{5.4cm}L{5.4cm}@{}}
\toprule
\rowcolor{primary!15}
 & \textbf{Symbolic (this course)} & \textbf{Learned (current wave)} \\
\midrule
\endhead
Where knowledge comes from & A person writes it down; the knowledge acquisition
bottleneck & Data; no bottleneck, and no way to correct one fact \\
Behaviour at the edge & Brittle: confidently wrong or silent & Degrades smoothly,
and gives no signal either \\
Guarantees & Real: a constraint solver proves no engineer is double-booked &
None \\
Explanation & A proof or a firing chain, checkable line by line & Text resembling
one, generated the same way as the answer \\
Ambiguous or novel input & Fails, often usefully & Handles it, sometimes by
inventing \\
\bottomrule
\caption{The two traditions, by their failure modes rather than their
successes.}\label{tab:neurosymbolic}
\end{longtable}

\term{Neuro-symbolic} systems combine them, and the combination that works
commercially is unglamorous: use the model for what is genuinely hard to specify ---
turning messy natural language into a structured request --- and a symbolic
component for what must be \emph{right}.

\begin{conceptbox}
Concretely, in the terms of this handout: a language model reads the ticket subject
and produces a structured request (\chref{nlp}'s stage~5, done better than any
grammar manages); \chref{csp}'s constraint solver then produces a schedule that
provably violates nothing; \chref{expert}'s rule engine applies the escalation
policy and emits the chain that an auditor will read.

Each component does the thing it is good at, and --- this is the part worth
remembering --- \emph{the guarantees come from the symbolic components}. A system
that needs to demonstrate that no engineer was double-booked, or to show exactly why
a ticket was escalated, gets that from Chapters~7 and~13 and cannot get it from the
model.

That is the concrete answer to ``why did I learn this?'' It is not that symbolic AI
will displace the current wave. It is that the parts of a deployed system which have
to be \emph{correct} and \emph{accountable} are still built the way this course
builds them, and the people who can build both are rare.
\end{conceptbox}

\section{Reading a Claim About AI}

The most transferable thing this chapter can leave you with is a procedure for
claims, because you will meet many and most will be overstated.

\begin{enumerate}[leftmargin=2.2em, itemsep=3pt, topsep=4pt]
  \item \textbf{Apply the five-point checklist} (\chref{introduction}). Goal or
        procedure? Does it consider alternatives? Is there an inspectable
        representation? Does it adapt? Can it explain? The answers are usually more
        interesting than the marketing.
  \item \textbf{Fill in the six columns} (\chref{classic}), and do not skip
        \emph{knowledge source} and \emph{failure mode}. Ask specifically: when it
        is wrong, how would anyone find out?
  \item \textbf{Ask what the demonstration excluded.} Student worked on a fragment
        of English. Every impressive demonstration is a microworld until shown
        otherwise.
  \item \textbf{Distinguish a claim about capability from a claim about reach.}
        \chref{introduction}'s history is two winters caused by exactly this
        conflation, and \chref{learning} showed a correct result about perceptrons
        producing a wrong conclusion about neural networks.
  \item \textbf{Ask what would falsify it.} A claim with no failing case is not a
        strong claim; it is an unfalsifiable one.
\end{enumerate}

\begin{worked}[Assessing a code-generation assistant]
A vendor claims their tool ``understands your codebase and writes correct code''.

\smallskip
\textbf{Checklist.} \emph{Goal:} no --- it is given a procedure (complete this
function), not an objective. \emph{Search:} sometimes; several candidates may be
generated and ranked. \emph{Knowledge:} partially --- if it uses the language
server's symbol table, that part is inspectable; the model's parameters are not.
\emph{Adaptation:} not from your corrections, generally. \emph{Explanation:} it will
produce a rationale on request, and that rationale is generated, not recorded. Score
around two, with the two most auditable tests failing.

\smallskip
\textbf{Six columns.} \emph{Failure mode} is the decisive one: it fails by producing
fluent, plausible, wrong code. Compare a type checker, which fails by refusing to
compile. The first needs an external detector; the second is one.

\smallskip
\textbf{What the demonstration excluded.} Almost certainly: a codebase with
unusual conventions, a language with little public code, and any case where the
right answer is ``this function should not exist''.

\smallskip
\textbf{Capability or reach?} The capability claim --- it writes usable code often
--- is well supported. The reach claim --- it \emph{understands} the codebase ---
is the \chref{nlp} Eliza effect restated, and the test is not whether the output
reads as understanding but whether there is a representation to point at.

\smallskip
\textbf{What would falsify it?} ``Correct code'' is testable and the vendor has not
offered a test. Ask for the acceptance criterion; the answer tells you whether they
have one.

\smallskip
\textbf{The useful conclusion.} Not ``it is good'' or ``it is hype'', but an
operational rule of the kind \chref{classic} produced: \emph{this tool's output goes
through the same review and test gate as a junior developer's, for the same reason,
and any workflow that bypasses that gate has removed the only failure detector in
the system}. That is a decision a team can act on, and it came from two checklists
rather than from an opinion.
\end{worked}

\begin{pitfall}
\begin{tight}
  \item \textbf{Treating generated text about reasoning as a trace of reasoning.}
        It is produced by the same process as the answer. \chref{expert}'s firing
        chain is a record; this is not.
  \item \textbf{Expecting hallucination to be patched away.} It is the objective
        working as specified. Mitigations reduce it; nothing removes it.
  \item \textbf{Assuming the agent loop is the model.} The loop, the tools, the
        budget and the termination test are ordinary engineering, and most failures
        live there.
  \item \textbf{Believing retrieval grounds an answer.} It puts text in the prompt.
        Nothing forces the answer to respect it.
  \item \textbf{Concluding symbolic methods are obsolete.} The guarantees and the
        audit trails in a deployed system come from them, and nothing else provides
        either.
  \item \textbf{Judging by the transcript.} That is the Eliza effect, and it worked
        on people who had read Eliza's source code.
\end{tight}
\end{pitfall}

%---------------------------------------------------------------------
\begin{lab}[A Tool-Using Agent, Offline]
The loop of \dsref{agent-tools}, with \chref{expert}'s rule engine where a language
model would sit. Deterministic, offline, and fully traceable --- which is what makes
it possible to see what the loop contributes.
\begin{lstlisting}[language=Python]
"""Chapter 20 lab -- a tool-using agent with no model in it.

The architecture of a modern agent: a goal, a set of tools, a loop
that decides what to call next, a step budget and a termination test.
The decision procedure here is Chapter 13's rule engine rather than a
language model, so every run is reproducible and every step has a
reason. No network, no API, no randomness.
"""

# ---- the world the tools read --------------------------------------
TASKS = {
    "auth-refresh":  {"points": 5, "depends": ["schema-change"]},
    "schema-change": {"points": 3, "depends": []},
    "ui-polish":     {"points": 2, "depends": []},
}
SPRINT_CAPACITY = 8


# ---- three tools. Each is cheap, local and honest about failure ----
def tool_estimate(task):
    t = TASKS.get(task)
    return t["points"] if t else None


def tool_dependencies(task):
    t = TASKS.get(task)
    return list(t["depends"]) if t else None


def tool_capacity(_):
    return SPRINT_CAPACITY


TOOLS = {
    "estimate": tool_estimate,
    "dependencies": tool_dependencies,
    "capacity": tool_capacity,
}


# ---- the decision procedure: Chapter 13's rule engine ---------------
# (name, [conditions on what we know], action)
# an action is ("call", tool, argument) or ("answer",)
POLICY = [
    ("need-deps",
     lambda k: "deps" not in k,
     lambda k: ("call", "dependencies", k["goal"])),
    # a tool can fail. A policy with no rule for that is the commonest
    # defect in deployed agents -- see the second experiment below.
    ("unknown-task",
     lambda k: k.get("deps", []) is None,
     lambda k: ("give-up",)),
    ("need-dep-estimates",
     lambda k: k.get("deps") and any(d not in k["points"]
                                     for d in k["deps"]),
     lambda k: ("call", "estimate",
                next(d for d in k["deps"] if d not in k["points"]))),
    ("need-own-estimate",
     lambda k: "deps" in k and k["goal"] not in k["points"],
     lambda k: ("call", "estimate", k["goal"])),
    ("need-capacity",
     lambda k: "capacity" not in k,
     lambda k: ("call", "capacity", None)),
    ("have-everything",
     lambda k: True,
     lambda k: ("answer",)),
]

BUDGET = 8


def agent(goal, policy=None, trace=True):
    """The loop of Figure 20.1. Keeps state, holds a goal, selects
    actions -- Chapter 2's model-based, goal-based agent."""
    policy = POLICY if policy is None else policy
    known = {"goal": goal, "points": {}}
    steps = []
    for step in range(1, BUDGET + 1):
        # MATCH and RESOLVE: first rule whose conditions hold
        action, name = None, None
        for (nm, cond, act) in policy:
            if cond(known):
                action, name = act(known), nm
                break
        if action is None:
            return None, steps, "no rule applies"
        if action[0] == "answer":
            steps.append((step, name, "answer", None))
            break
        if action[0] == "give-up":
            steps.append((step, name, "give up", None))
            return None, steps, "unknown task: %s" % goal
        _, tool, arg = action
        result = TOOLS[tool](arg)
        steps.append((step, name, "%s(%s)" % (tool, arg), result))
        # ACT: fold the observation into what we know
        if tool == "dependencies":
            known["deps"] = result
        elif tool == "estimate":
            known["points"][arg] = result
        elif tool == "capacity":
            known["capacity"] = result
    else:
        return None, steps, "budget exhausted"

    total = sum(known["points"].values())
    fits = total <= known["capacity"]
    answer = ("%s plus its dependencies is %d points against a "
              "capacity of %d: %s"
              % (goal, total, known["capacity"],
                 "it fits" if fits else "it does NOT fit"))
    return answer, steps, None


if __name__ == "__main__":
    print("goal: can we take 'auth-refresh' into the next sprint?")
    print()
    answer, steps, err = agent("auth-refresh")
    print("step  rule                 action                   result")
    print("-" * 66)
    for (n, rule, action, result) in steps:
        print("%4d  %-20s %-24s %s" % (n, rule, action, result))
    print()
    print("ANSWER:", answer)
    assert err is None
    # 5 points for auth-refresh, 3 for the dependency it pulled in,
    # against a capacity of 8: it fits exactly, with nothing spare
    assert "it fits" in answer

    # every step has a named reason, and the run is reproducible
    a2, s2, _ = agent("auth-refresh")
    assert [x[1] for x in s2] == [x[1] for x in steps]
    print()
    print("Run it again and the trace is identical: the decision")
    print("procedure is a rule base, so the agent is reproducible")
    print("and every step names the rule that caused it.")

    # --- a tool that fails, handled --------------------------------
    print()
    print("=== when a tool returns nothing ===")
    a3, s3, err3 = agent("does-not-exist")
    print("   goal 'does-not-exist' ->", err3)
    print("   rule that caught it:", s3[-1][1])
    assert err3.startswith("unknown task")
    print("   One rule, 'unknown-task', turns a failed tool call into")
    print("   an honest answer rather than a crash.")

    # --- and what happens without it -------------------------------
    print()
    print("=== the same agent with that rule removed ===")
    careless = [r for r in POLICY if r[0] != "unknown-task"]
    try:
        a4, s4, err4 = agent("does-not-exist", policy=careless)
        print("   result:", err4 or a4)
    except TypeError as exc:
        print("   TypeError: %s" % exc)
        print("   The tool returned None, no rule noticed, and the")
        print("   None travelled all the way to the arithmetic at the")
        print("   end. This is the commonest failure in deployed agent")
        print("   systems, and note where it lives: in the POLICY, not")
        print("   in the model at the centre.")

    # the step budget is the other guard
    print()
    print("=== what the step budget is for ===")
    spinner = [("always-ask", lambda k: True,
                lambda k: ("call", "capacity", None))]
    a5, s5, err5 = agent("auth-refresh", policy=spinner)
    print("   a policy that never terminates ->", err5)
    print("   it ran %d steps and was stopped by the budget." % len(s5))
    assert err5 == "budget exhausted" and len(s5) == BUDGET

    # --- what is symbolic and what would be learned ----------------
    print()
    print("=== which part would a language model replace? ===")
    print("   the LOOP            -- no, ordinary engineering")
    print("   the TOOLS           -- no, they read a database")
    print("   the STEP BUDGET     -- no, and it is load-bearing")
    print("   the TERMINATION     -- no")
    print("   the DECISION RULE   -- YES, this is the part")
    print()
    print("   Five of the six components are unchanged. Most of what")
    print("   makes an agent work or fail lives in the five.")
\end{lstlisting}
Then delete the \texttt{BUDGET} check and re-run the second example. The agent loops
forever on a goal it cannot resolve, because no rule consumes a \texttt{None}. That
is the commonest failure in deployed agent systems, and it has nothing to do with
the model at the centre.
\end{lab}

%---------------------------------------------------------------------
\begin{checkpoint}
\begin{tightnum}
  \item A language model produces a step-by-step explanation of its answer. State
        why this is not the same artefact as \chref{expert}'s firing chain, using
        the distinction between $\models$ and $\vdash$ from \chref{propositional}.
  \item The lab's agent has six components and a language model would replace one.
        Name the other five and say which one, removed, causes the commonest
        deployed failure.
  \item A vendor demonstrates a system that answers questions about your codebase
        fluently and correctly in the demonstration. Give the two questions from
        this chapter's procedure that would tell you most, and say why.
\end{tightnum}
\end{checkpoint}

%---------------------------------------------------------------------
\begin{chaptersummary}
\begin{tight}
  \item A \term{large language model} predicts the next \term{token}. It is
        \chref{nlp}'s n-gram idea with a much larger window learned rather than
        counted, which addresses exactly the limitation counting could not.
  \item There is no representation to inspect, no recorded derivation, and
        \term{hallucination} is the objective working as specified rather than a
        defect. Its failure mode is Eliza's: none detectable by the system.
  \item A \term{tool-using agent} is \chref{agents}' model-based, goal-based agent.
        Only the decision procedure at its centre is new; the loop, tools, budget
        and termination test are ordinary engineering, and most failures live
        there.
  \item \term{Retrieval} grounds an answer by putting text in the prompt. Nothing
        forces the answer to respect it.
  \item The two traditions fail in opposite directions. \term{Neuro-symbolic}
        systems use the model where specification is hard and a symbolic component
        where the answer must be right --- and \emph{the guarantees and the audit
        trails come from the symbolic components}.
  \item To assess a claim: apply \chref{introduction}'s checklist, fill
        \chref{classic}'s six columns, ask what the demonstration excluded,
        separate capability from reach, and ask what would falsify it.
  \item Judging by the transcript is the Eliza effect. It worked on people who had
        read the source code.
\end{tight}
\end{chaptersummary}

\begin{reviewq}
  \item \textbf{(MCQ)} A large language model's training objective is to:
        (a)~answer questions correctly (b)~predict the next token
        (c)~maximise expected utility (d)~satisfy constraints.
  \item \textbf{(MCQ)} Hallucination is best described as: (a)~a bug to be patched
        (b)~the training objective working as specified (c)~a retrieval failure
        (d)~a tokenisation error.
  \item \textbf{(MCQ)} In a tool-using agent, a language model typically replaces:
        (a)~the tools (b)~the loop (c)~the decision about what to do next
        (d)~the step budget.
  \item \textbf{(MCQ)} In a neuro-symbolic system, the guarantees come from:
        (a)~the model (b)~the symbolic component (c)~retrieval (d)~the prompt.
  \item \textbf{(Short)} Explain how a large language model addresses the n-gram
        limitation of \chref{nlp}, and what it gives up that a grammar provided.
  \item \textbf{(Short)} Why is a generated explanation not a trace? Answer using
        the $\models$ / $\vdash$ distinction.
  \item \textbf{(Short)} Give a delivery task you would hand to a model and one you
        would hand to a constraint solver, and justify each in one sentence.
  \item \textbf{(Applied)} Take a current AI product you have used. Fill in
        \chref{classic}'s six columns for it, filling what you can from outside,
        and state the one operational rule your failure-mode analysis produces.
\end{reviewq}
